% TOP_PROBLEM: {"id": 1860, "title": "Shanks Chains of Length 7", "category_id": 1, "category": {"id": 1, "name": "number_theory", "display_name": "Number Theory", "description": "Properties of integers, prime numbers, Diophantine equations.", "slug": "number-theory", "order_index": 1, "created_at": "2026-07-31T15:26:25.670Z"}, "status": "open"}
% ATTEMPT_STATUS: unresolved
% Research attempt by GPT_6_Astra_Ultra, 1 October 2026.
% Canonical UnsolvedMath ID; original statements and sources preserved.
% FIRST_POSTED: 2026-10-01
% Imported from: unsolvedmath_additional_500.tex; UnsolvedMath ID 1860
% !TeX program = lualatex
% Compile twice: lualatex 1860.tex
% Self-contained: no external bibliography, images, or \input files.
% Numeric section labels and visible numbers are original UnsolvedMath IDs.
\documentclass[11pt,a4paper]{article}
\usepackage[margin=24mm,headheight=14pt]{geometry}
\usepackage{fontspec}
\setmainfont{Latin Modern Roman}
\setsansfont{Latin Modern Sans}
\setmonofont{Latin Modern Mono}
\usepackage{amsmath,amssymb,mathtools}
\usepackage{unicode-math}
\setmathfont{Latin Modern Math}
\usepackage{microtype}
\usepackage{enumitem,xurl,needspace}
\usepackage[dvipsnames]{xcolor}
\usepackage{hyperref}
\hypersetup{unicode=true,colorlinks=true,linkcolor=MidnightBlue,urlcolor=MidnightBlue,
 pdftitle={UnsolvedMath Problem 1860},pdfauthor={GPT 6 Astra Ultra},bookmarksnumbered=true}
\usepackage{bookmark,tocloft,titlesec,fancyhdr}
\setlength{\cftsecnumwidth}{4.2em}
\cftsetpnumwidth{2.5em}
\cftsetrmarg{4em}
\titleformat{\section}{\Large\bfseries\raggedright}{\thesection}{0.7em}{}
\pagestyle{fancy}\fancyhf{}
\fancyhead[L]{\small\sffamily GPT 6 Astra Ultra research attempt}
\fancyhead[R]{\small Original catalogue IDs}
\fancyfoot[C]{\thepage}
\setlength{\parindent}{0pt}
\setlength{\parskip}{5pt plus 1pt}
\setlength{\emergencystretch}{3em}
\setcounter{tocdepth}{1}\setcounter{secnumdepth}{1}
\setlist[itemize]{leftmargin=3em,itemsep=4pt,topsep=4pt,parsep=0pt}
\allowdisplaybreaks
\newcommand{\N}{\mathbb{N}}
\newcommand{\Z}{\mathbb{Z}}
\newcommand{\Q}{\mathbb{Q}}
\newcommand{\R}{\mathbb{R}}
\newcommand{\C}{\mathbb{C}}
\newcommand{\F}{\mathbb{F}}
\newcommand{\A}{\mathbb{A}}
\newcommand{\PP}{\mathbb{P}}
\newcommand{\E}{\mathbb{E}}
\newcommand{\eps}{\varepsilon}
\newcommand{\dd}{\,\mathrm d}
\newcommand{\sref}[2]{\hyperlink{source:#1:#2}{[S#2]}}
\newif\ifshowcatalogue
\showcataloguetrue
\newcommand{\cataloguescope}[1]{\ifshowcatalogue
 \paragraph{Statement recorded in the supplied source.}
 #1\fi}
\begin{document}
% UnsolvedMath ID 1860; source catalogue code GUY-A7b

\setcounter{section}{1859}

\section{A Shanks Prime Chain of Length Seven}\label{1860}

{\small\sffamily UnsolvedMath ID 1860 · Number Theory}

\subsection{Definitions and mathematical statement}

\paragraph{Shared notation.}Unless a section says otherwise, $\N=\{1,2,\ldots\}$,
$\Z$ is the ring of integers, $\Q,\R,\C$ are the rational, real, and complex
fields, $[N]=\{1,\ldots,\lfloor N\rfloor\}$, and $\log$ is the natural
logarithm. For real functions of a parameter tending to infinity,
$f\ll g$ and $f=O(g)$ mean $|f|\leq Cg$ eventually for some fixed $C$;
$f\gg g$ reverses this comparison; $f=o(g)$ means $f/g\to0$;
$f\sim g$ means $f/g\to1$; and $f\asymp g$ means both order bounds.
A subscript on $O$ or $\ll$ specifies allowed constant dependence.
The expression $n^{o(1)}$ in an upper bound means growth smaller than
$n^\epsilon$ for each fixed $\epsilon>0$. Local definitions govern any
exception, finite-field convention, or use of the same symbol for another
object. An asymptotic claim is not automatically an effective algorithm.



A chain of length seven is a list of seven primes $p_1,\ldots,p_7$ with $p_{i+1}=4p_i^2-17$ for $1\leq i\leq6$. The first prime is not prescribed. All seven terms must be positive primes; changing the recurrence or counting seven transitions would define a different search problem. \sref{1860}{1} \sref{1860}{2}

\paragraph{Statement recorded in the supplied source.}
Are there any Shanks chains of length 7 with $p_{i+1} = 4p_i^2 - 17$? \sref{1860}{1}

\subsection{Short English statement}

Can seven consecutive terms of the recurrence four times the square minus seventeen all be prime?

\subsection{Research attempt}

\paragraph{A catalogue-status discrepancy must be resolved first.}
The supplied statement labels existence of a seven-term chain as open,
but its own reference [S2] reports a length-seven example found by
Paul Leyland on 22 November 2000, starting at $176996212597$, with
a $758$-digit seventh term. This is a reported prior resolution of the
literal existence question, not a new conjectural candidate invented
here. The cited primary announcement [E1] was sought on 1 October 2026
but could not be retrieved. Thus this notebook performs a mathematical
and computational audit of the reported example and records the exact
limit of independent certification, rather than falsely declaring the
literature question still open or claiming a new solution.

\paragraph{Exact reconstruction of the seven reported terms.}
Let $F(t)=4t^2-17$ and set
\[
 p_1=176996212597,\qquad p_j=F^{j-1}(p_1)\quad(2\le j\le7).
\]
This specifies seven integers with six transitions, precisely matching
the catalogue convention. The script executes all six integer
recurrences and writes every term in full to the result file. Their
decimal digit counts are
\[
                         12,24,47,95,189,379,758.
\]
The last count agrees with the catalogue reference's report. Exact
agreement of recurrence and size is useful evidence of identifying the
reported chain, but is not a primality certificate for the terms.

\paragraph{Growth and positivity are elementary.}
For every integer $t\ge3$,
\[
 F(t)-t=4t^2-t-17>0,
\]
so a chain starting above two remains positive and strictly increases.
Starting at the prime two fails immediately, because $F(2)=-1$.
For $t\ge3$, the bounds $2t^2<F(t)<4t^2$ hold, giving
\[
 2\log t+\log2<\log F(t)<2\log t+\log4.
\]
Thus bit length roughly doubles at each transition. This explains why
checking the first few primes is qualitatively different from proving
the final hundreds-of-digits term prime, despite the small number of
transitions.

\paragraph{An exact congruence sieve for potential starts.}
For a fixed prime $\ell$, reduce the map $F$ modulo $\ell$.
A start residue can be rejected if one of its first seven iterates
is zero modulo $\ell$, provided the corresponding integer term is
larger than $\ell$. Equality to $\ell$ is a prime exception and must
not be discarded by an indiscriminate divisibility test.

For example modulo thirteen,
$F(t)\equiv4(t^2-1)$. The residues $\pm1$ reach zero in one step,
and $\pm2$ reach $-1$ and then zero. Direct enumeration shows that
the residues surviving seven nonzero terms are
\[
                          3,4,5,6,7,8,9,10\pmod{13}.
\]
For starts larger than thirteen, every actual term is larger than
thirteen by growth, so this is a sound necessary condition for a
seven-prime chain. For smaller starts the equality exception must be
tested using the actual integer term.

For any finite list of primes one can intersect the allowed residue
conditions using the Chinese remainder theorem. This is a finite
search filter, not a proof of simultaneous primality. A residue
surviving the small-prime filters can still yield a composite term
whose factors exceed all sieve primes.

\paragraph{What was and was not certified locally.}
The script proves $p_1$ prime by trial division through its integer
square root. It reconstructs all later terms exactly, and separately
enumerates seven-step nonzero residue classes for small moduli. It
does not claim a deterministic primality certificate for $p_2$ through
$p_7$. No probable-prime test is silently promoted to a primality
proof; indeed none is needed for the stated local results. The result
file explicitly marks those six primalities as not certified here.

A complete independent verification could use a checked ECPP
certificate, or another rigorous primality certificate, for each
later term. Merely testing many Fermat bases would not close that gap.
Likewise a chain of probable primes would meet a different, weaker
search specification than the seven actual primes required here.

\paragraph{Artifacts and literature boundary.}
The exact terms and digit counts are stored in
\texttt{results\_second.json}, generated by
\texttt{checks\_second.py}, under
\path{_Local/Erdos_problems/UnsolvedMath/attempt_audit_2026_10_01/number_theory/}.
The source report supplies a named prior example, while the independent
arithmetic verifies recurrence, positivity, term counts, and the first
primality. These are distinct evidence levels and are kept separate.
No claim to have contacted the original author or obtained unavailable
certificates is made.

\paragraph{Outcome.}
\textbf{Prior existence reported; local certification UNRESOLVED.}
The literal catalogue question is reported answered by the existing
source. This attempt does not independently certify all seven primes
and therefore retains an unresolved attempt marker, referring to that
audit gap rather than asserting a fresh open status for the published
existence question. A reviewer should update catalogue status only
after evaluating the original announcement or an equivalent certificate.

\subsection{Sources}

{\small

\begin{itemize}

\item[\hypertarget{source:1860:1}{S1}] ULAM.AI, \emph{UnsolvedMath}, supplied \texttt{problems.json}, record 1860 (GUY-A7b), statement and background. The embedded literature-review date is 2026-08-17; that is the archive’s review date, not a fresh certification. Dataset landing page: \url{https://huggingface.co/datasets/ulamai/UnsolvedMath}.

\item[\hypertarget{source:1860:2}{S2}] Shanks Chain, Wolfram MathWorld. \url{https://mathworld.wolfram.com/ShanksChain.html}. Bibliographic information imported from the supplied record.


\item[E1] Paul Leyland, \emph{A Shanks (4,17) Chain of Length 7},
NMBRTHRY announcement, 22 November 2000, as linked by [S2].
\url{https://listserv.nodak.edu/cgi-bin/wa.exe?A2=ind0011&L=nmbrthry&P=2651}.
Primary-source retrieval attempted 1 October 2026 but unavailable.
Its content is not claimed to have been independently inspected.
\item[E2] Eric W. Weisstein, \emph{Shanks Chain}, cited catalogue
reference, inspected 1 October 2026 for the reported start and length.
\url{https://mathworld.wolfram.com/ShanksChain.html}.
This is the secondary lead to [E1], not an independently checked
primality certificate for the final terms.

\end{itemize}

}

\label{end:1860}
\end{document}
